\documentclass[12pt]{article}
\usepackage[utf8]{inputenc}
\usepackage[T1]{fontenc}
\usepackage{hyperref}
\title{Samaritan Aramaic Language Resources for K-12 Curriculum Development}
\author{Compiled by Claude Sonnet 5.5 (claude-sonnet-5-5) for Andr\'as Kornai}
\date{}
\begin{document}
\maketitle

\section{Introduction}
Samaritan Aramaic (ISO 639-3 \texttt{sam}) is the literary and liturgical Aramaic of the Samaritan community in Palestine and Israel. It has no community of first-language speakers and is not known to be taught in schools. The resources found are scholarly.

\section{Primary Educational Resources}
No school primers or textbooks were found.

\section{Linguistic and Reference Materials}
\begin{itemize}
\item Abraham Tal, \emph{A Dictionary of Samaritan Aramaic} (Brill, 2000, 2 vols.). It covers the Aramaic dialect of Samaritan texts from the 4th to the 11th century, with entries in Hebrew script and English and Hebrew glosses. Publisher listing: \url{https://www.degruyterbrill.com/document/isbn/9789004294158/html}. Library record: \url{https://vufind.lib.tsu.ru/Record/libtsu.1013862/Description}.
\item \emph{Brevis Linguae Samaritanae} (grammar, chrestomathy and glossary; an old work, found only as a facsimile reprint listing): \url{https://buecher.de/artikel/buch/brevis-linguae-samaritanae/26040247}
\end{itemize}

\section{Religious and Community Resources}
The language is tied to Samaritan liturgy and Targum literature. No community-run teaching resources were found in the searches.

\section{Digital and Audio Resources}
No audio or digital resources were found.

\section{Curriculum Development Context}
Little exists beyond scholarly dictionaries and an old grammar. Any K-12 use would have to rely on liturgical texts and adaptation of university-level materials.

\end{document}